\documentclass[a4paper,12pt]{article}

\usepackage[utf8]{inputenc}
\usepackage{amsmath}
\usepackage{amssymb}
\usepackage{geometry}
\usepackage{graphicx}
\usepackage{float}
\usepackage{siunitx}
\usepackage{listings}
\usepackage{hyperref}
\usepackage{booktabs}
\usepackage{amsfonts}

\geometry{margin=2.5cm}

\sisetup{exponent-product = \cdot, output-product = \cdot}
\lstset{
    breaklines=true,
    breakatwhitespace=true,
    columns=fullflexible
}

\title{IoT Homework -- Exercise 2}
\author{Angelica Ortiz 11191225, Daniel Nägele 11180610}
\date{May 2026}

\begin{document}

\maketitle
\thispagestyle{empty}
\newpage

\tableofcontents
\newpage
\section{Probability Mass Function}
\subsection{Logic}
The assumed main use case for the localization of these dentures is the loss and the prevention of loss of said dentures. It makes no sense to use them to localize patients, as not all patients are wearing dentures in general, dentures are not worn at all times, and the fact that the human body will absorb signals. Especially the \SI{2.4}{\giga\hertz} band used by BLE resonates heavily with water molecules and thus the human body. This will lead to absorption, scattering and reflection, and result in severe skewing of RSSI measurements.\\
This leads us to consider the following three scenarios:
\vspace{5pt}

\noindent\textbf{In-Mouth:} While the denture is worn, we want to put the device to deep sleep and not transmit any beacon signals. Only a signal telling the system that the denture is worn would be useful.

\vspace{5pt}

\noindent\textbf{Out of Mouth, Stationary:} When the denture is resting on a nightstand or in a glass, only a few initial beacons are necessary to estimate the location. As long as the dentures remain stationary, no further beacons need to be emitted.

\vspace{5pt}

\noindent\textbf{Out of Mouth, Moving:} If the denture is not worn but still moving, the beacons need to be emitted continuously and often. Such a case would be if the denture is accidentally swept into laundry or the trash and the loss of device needs to be prevented. Geo-fencing could be implemented in the sensing side of the system to also automatically alert caregivers.

\subsection{Sensing}
To detect our scenarios of interest, we need two additional sensors.
\begin{itemize}
	\item Temperature Sensor: Used to detect if the denture is inside the mouth by measuring the body's heat (e.g. temperature t \textgreater{} \SI{34}{\degreeCelsius})
	\item Acceleration Sensor: Used to detect motion if the denture is outside the mouth.
\end{itemize}

\subsection{Pseudocode}

\begin{lstlisting}
resting_transmits = 0

setup() {
  init_temp_sensor()
  init_accelerometer(Wake_on_Motion_Interrupt) // send interrupt if motion is detected
}

loop() {
  temperature = read_temperature()
  motion_detected = read_last_motion_state()

  if temperature > 34.0:
    resting_transmits = 0
    disable_ble_beacons()
    ignore_interrupts(Wake_on_Motion_Interrupt) // while in-mouth, ignore unnecessary interrupt-based wakeups
    sleep(5 minutes) // acts as return statement, no further code is executed

  enable_interrupts(Wake_on_Motion_Interrupt)
  if motion_detected:
    resting_transmits = 0
    enable_ble_beacons()
    set_ble_beacon_frequency(1 Hz)
    sleep(1 second)

  // we have transmitted our beacon 5 times and will not emit any further beacons until movement is detected
  if resting_transmits is 5:
    disable_ble_beacons()
    resting_transmits = 0
    sleep(1 day) // no need to wake up before motion is detected, only 1 day as a fail-safe
  else:
    enable_ble_beacons()
    set_ble_beacon_frequency(1/60 Hz) // transmit beacon every minute
    resting_transmits = resting_transmits + 1
    sleep(1 minute)
}

\end{lstlisting}
\end{document}
